The Presentation Layer implements all user-visible behavior described at a high level in the SRS: authentication, dashboards, plant health display, historical views, scheduling, mesh-oriented navigation, and Summer~2026 \textbf{edge device} controls (Devices list, tray Raspberry~Pi panel, Take Picture). Each subsystem below corresponds to a coherent screen family or cross-cutting UI concern.

\subsection{Authentication and session UI subsystem}

This subsystem collects credentials or OAuth flows, exchanges tokens for server sessions, and guards entry to protected routes. It satisfies SRS \textbf{user authentication} expectations at the UX level.

\begin{figure}[H]
\centering
\begin{tikzpicture}[font=\scriptsize, box/.style={draw, rounded corners, align=center, minimum width=2.2cm, minimum height=0.7cm}, arr/.style={-{Stealth}, thick}]
\node[box, fill=leafgreen!15] (login) {Login /\\register};
\node[box, right=1.0cm of login, fill=blue!8] (api) {Application\\APIs};
\draw[arr] (login) -- node[above] {HTTPS} (api);
\end{tikzpicture}
\caption{Authentication UI: user actions become authenticated API calls.}
\label{fig:pres-auth}
\end{figure}

\subsubsection{Assumptions}
\begin{itemize}
  \item Users have network connectivity for initial sign-in.
  \item Identity provider availability matches organizational SLA.
\end{itemize}

\subsubsection{Responsibilities}
\begin{itemize}
  \item Render sign-in, registration, and sign-out affordances.
  \item Hold short-lived client tokens only as required by the chosen identity integration.
  \item Redirect unauthenticated visitors away from protected console routes.
\end{itemize}

\subsubsection{Subsystem interfaces}
\begin{table}[H]
\centering
\small
\begin{tabular}{|p{0.9cm}|p{3.4cm}|p{3.0cm}|p{3.5cm}|}
\hline
\textbf{ID} & \textbf{Description} & \textbf{Inputs} & \textbf{Outputs} \\
\hline
P-A1 & User $\rightarrow$ Auth UI & Credentials / OAuth & Sign-in intent \\
\hline
P-A2 & Auth UI $\rightarrow$ Application & Session exchange requests & HTTP-only session \\
\hline
\end{tabular}
\caption{Authentication and session UI subsystem interfaces}
\end{table}

\subsection{Dashboard and overview subsystem}

The \textbf{dashboard} aggregates tray summaries, recent captures or predictions where available, and entry points into deeper views. It supports SRS \textbf{dashboard interface} and high-level plant status visibility.

\begin{figure}[H]
\centering
\begin{tikzpicture}[font=\scriptsize, box/.style={draw, rounded corners, align=center, minimum width=2.2cm, minimum height=0.7cm}, arr/.style={-{Stealth}, thick}]
\node[box, fill=leafgreen!12] (dash) {Dashboard\\widgets};
\node[box, right=1.0cm of dash, fill=blue!8] (api) {Application};
\draw[arr] (dash) -- node[above] {fetch} (api);
\end{tikzpicture}
\caption{Dashboard subsystem requests summary aggregates from the application tier.}
\label{fig:pres-dash}
\end{figure}

\subsubsection{Assumptions}
\begin{itemize}
  \item Summary data volume is bounded for interactive load times.
  \item Chart libraries render only on supported browsers.
\end{itemize}

\subsubsection{Responsibilities}
\begin{itemize}
  \item Display tray list and key health indicators.
  \item Link to tray detail, plants, and schedules.
  \item Surface empty and error states clearly.
\end{itemize}

\subsubsection{Subsystem interfaces}
\begin{table}[H]
\centering
\small
\begin{tabular}{|p{0.9cm}|p{3.4cm}|p{3.0cm}|p{3.5cm}|}
\hline
\textbf{ID} & \textbf{Description} & \textbf{Inputs} & \textbf{Outputs} \\
\hline
P-D1 & Dashboard $\rightarrow$ Application & Filter context & Summary DTOs \\
\hline
P-D2 & Application $\rightarrow$ Dashboard & JSON payloads & Rendered charts / lists \\
\hline
\end{tabular}
\caption{Dashboard and overview subsystem interfaces}
\end{table}

\subsection{Tray monitoring console subsystem}

The \textbf{tray console} shows the active tray: latest imagery, monitoring trends, plant thumbnails, events, optional tray-level vision uploads, and the embedded \textbf{Raspberry~Pi panel}. It aligns with \textbf{automatic image capture} presentation and tray-scoped monitoring in the SRS.

\begin{figure}[H]
\centering
\begin{tikzpicture}[font=\scriptsize, box/.style={draw, rounded corners, align=center, minimum width=2.0cm, minimum height=0.65cm}, arr/.style={-{Stealth}, thick}]
\node[box, fill=leafgreen!14] (ui) {Tray UI};
\node[box, right=0.9cm of ui, fill=blue!8] (api) {API};
\node[box, right=0.9cm of api, fill=orange!10, dashed] (cv) {Vision\\(opt.)};
\draw[arr] (ui) -- (api);
\draw[arr, dashed] (api) -- (cv);
\end{tikzpicture}
\caption{Tray console: drives reads and optional vision requests through the application layer.}
\label{fig:pres-tray}
\end{figure}

\subsubsection{Assumptions}
\begin{itemize}
  \item Each tray identifier maps to authorized user scope.
  \item Large images may be progressively loaded or thumbnailed.
  \item A tray may be linked to zero or one \texttt{edge\_devices} row.
\end{itemize}

\subsubsection{Responsibilities}
\begin{itemize}
  \item Render tray imagery, charts, and plant rows.
  \item Submit tray-scoped uploads or vision jobs as permitted.
  \item Host the Raspberry~Pi panel affordances (link device, Take Picture, Klipper/streamer URL).
  \item Poll or refresh data according to UX policy.
\end{itemize}

\subsubsection{Subsystem interfaces}
\begin{table}[H]
\centering
\small
\begin{tabular}{|p{0.9cm}|p{3.4cm}|p{3.0cm}|p{3.5cm}|}
\hline
\textbf{ID} & \textbf{Description} & \textbf{Inputs} & \textbf{Outputs} \\
\hline
P-T1 & Tray UI $\rightarrow$ Application & Tray ID, actions & Tray detail DTO \\
\hline
P-T2 & Tray UI $\rightarrow$ Application & Image multipart & Vision job result \\
\hline
\end{tabular}
\caption{Tray monitoring console subsystem interfaces}
\end{table}

\subsection{Plant detail and health history subsystem}

\textbf{Plant detail} focuses on a single plant: editable metadata, photo gallery, health reports over time, and scoped monitoring logs. It directly supports SRS \textbf{plant health status} and \textbf{historical data display}.

\begin{figure}[H]
\centering
\begin{tikzpicture}[font=\scriptsize, box/.style={draw, rounded corners, align=center, minimum width=2.0cm, minimum height=0.65cm}, arr/.style={-{Stealth}, thick}]
\node[box, fill=leafgreen!14] (p) {Plant UI};
\node[box, right=1.0cm of p, fill=blue!8] (api) {Application};
\draw[arr] (p) -- node[above] {CRUD + history} (api);
\end{tikzpicture}
\caption{Plant detail subsystem: CRUD and reporting views.}
\label{fig:pres-plant}
\end{figure}

\subsubsection{Assumptions}
\begin{itemize}
  \item Users may edit only plants they own or that policy allows.
  \item Report history may be lengthy; pagination or lazy loading applies.
\end{itemize}

\subsubsection{Responsibilities}
\begin{itemize}
  \item Display health timeline and latest findings.
  \item Support photo upload and deletion with confirmation.
  \item Invoke ``add plant from photo'' flows when enabled.
\end{itemize}

\subsubsection{Subsystem interfaces}
\begin{table}[H]
\centering
\small
\begin{tabular}{|p{0.9cm}|p{3.4cm}|p{3.0cm}|p{3.5cm}|}
\hline
\textbf{ID} & \textbf{Description} & \textbf{Inputs} & \textbf{Outputs} \\
\hline
P-P1 & Plant UI $\rightarrow$ Application & Plant ID & Detail + reports \\
\hline
P-P2 & Plant UI $\rightarrow$ Application & Edits / uploads & Updated plant \\
\hline
\end{tabular}
\caption{Plant detail and health history subsystem interfaces}
\end{table}

\subsection{Mesh and topology subsystem}

\textbf{Mesh views} let users group trays for aggregated charts and cross-tray plant listings, supporting experimental layouts described in the SRS's organizational needs without mandating a particular hardware topology.

\begin{figure}[H]
\centering
\begin{tikzpicture}[font=\scriptsize, box/.style={draw, rounded corners, align=center, minimum width=2.0cm, minimum height=0.65cm}, arr/.style={-{Stealth}, thick}]
\node[box, fill=leafgreen!12] (m) {Mesh UI};
\node[box, right=1.0cm of m, fill=blue!8] (api) {Application};
\draw[arr] (m) -- node[above] {mesh CRUD} (api);
\end{tikzpicture}
\caption{Mesh subsystem interacts with mesh and member tray aggregates.}
\label{fig:pres-mesh}
\end{figure}

\subsubsection{Assumptions}
\begin{itemize}
  \item Mesh membership changes are infrequent relative to reads.
\end{itemize}

\subsubsection{Responsibilities}
\begin{itemize}
  \item Create and list meshes; show merged monitoring where available.
  \item Navigate from mesh to constituent trays and plants.
\end{itemize}

\subsubsection{Subsystem interfaces}
\begin{table}[H]
\centering
\small
\begin{tabular}{|p{0.9cm}|p{3.4cm}|p{3.0cm}|p{3.5cm}|}
\hline
\textbf{ID} & \textbf{Description} & \textbf{Inputs} & \textbf{Outputs} \\
\hline
P-M1 & Mesh UI $\rightarrow$ Application & Mesh ID / payload & Mesh DTO \\
\hline
\end{tabular}
\caption{Mesh and topology subsystem interfaces}
\end{table}

\subsection{Schedule management subsystem}

The \textbf{schedule editor} binds capture timing to trays or meshes, supporting SRS \textbf{automatic image capture} policies at the configuration level.

\begin{figure}[H]
\centering
\begin{tikzpicture}[font=\scriptsize, box/.style={draw, rounded corners, align=center, minimum width=2.0cm, minimum height=0.65cm}, arr/.style={-{Stealth}, thick}]
\node[box, fill=leafgreen!12] (s) {Schedule UI};
\node[box, right=1.0cm of s, fill=blue!8] (api) {Application};
\draw[arr] (s) -- node[above] {upsert} (api);
\end{tikzpicture}
\caption{Schedule subsystem persists policies via application APIs.}
\label{fig:pres-sch}
\end{figure}

\subsubsection{Assumptions}
\begin{itemize}
  \item Schedules are stored in a timezone-aware or UTC-normalized manner in the data layer.
\end{itemize}

\subsubsection{Responsibilities}
\begin{itemize}
  \item List and edit schedules with clear scope (tray vs.\ mesh).
  \item Support destination \texttt{raspberry-pi-edge} for edge capture policies.
  \item Validate user inputs (cadence, start/end) before submission.
\end{itemize}

\subsubsection{Subsystem interfaces}
\begin{table}[H]
\centering
\small
\begin{tabular}{|p{0.9cm}|p{3.4cm}|p{3.0cm}|p{3.5cm}|}
\hline
\textbf{ID} & \textbf{Description} & \textbf{Inputs} & \textbf{Outputs} \\
\hline
P-S1 & Schedule UI $\rightarrow$ Application & Schedule DTO & Persisted schedule \\
\hline
\end{tabular}
\caption{Schedule management subsystem interfaces}
\end{table}

\subsection{Devices UI subsystem}

The \textbf{Devices} navigation lists provisioned Raspberry~Pi edge devices for the signed-in owner: CPU serial, online/offline status, last heartbeat, and linked tray. It is the fleet-level companion to the per-tray Raspberry~Pi panel.

\begin{figure}[H]
\centering
\begin{tikzpicture}[font=\scriptsize, box/.style={draw, rounded corners, align=center, minimum width=2.2cm, minimum height=0.7cm}, arr/.style={-{Stealth}, thick}]
\node[box, fill=leafgreen!14] (dev) {Devices\\list UI};
\node[box, right=1.0cm of dev, fill=blue!8] (api) {/api/devices};
\draw[arr] (dev) -- node[above] {GET / actions} (api);
\end{tikzpicture}
\caption{Devices UI: operator fleet view over edge device APIs.}
\label{fig:pres-devices}
\end{figure}

\subsubsection{Assumptions}
\begin{itemize}
  \item Only devices owned by the authenticated account are visible.
  \item Heartbeat stale windows determine online presentation.
\end{itemize}

\subsubsection{Responsibilities}
\begin{itemize}
  \item Render device inventory and status.
  \item Navigate to linked trays when present.
  \item Surface revoke / rotate-key actions when exposed by product policy.
\end{itemize}

\subsubsection{Subsystem interfaces}
\begin{table}[H]
\centering
\small
\begin{tabular}{|p{0.9cm}|p{3.4cm}|p{3.0cm}|p{3.5cm}|}
\hline
\textbf{ID} & \textbf{Description} & \textbf{Inputs} & \textbf{Outputs} \\
\hline
P-DV1 & Devices UI $\rightarrow$ Application & Session context & Device list DTO \\
\hline
P-DV2 & Devices UI $\rightarrow$ Application & Device action payload & Updated device state \\
\hline
\end{tabular}
\caption{Devices UI subsystem interfaces}
\end{table}

\subsection{Tray Raspberry Pi panel and Take Picture subsystem}

The \textbf{tray Raspberry~Pi panel} embeds edge controls on tray detail: link/unlink device, edit and save \textbf{Klipper/streamer URL} (\texttt{klipper\_url}), generate poses from plant layout, and \textbf{Take Picture}. Take Picture posts \texttt{action=capture} to \texttt{/api/devices/\{id\}}; the UI primarily expects a queued-command acknowledgment for agent \texttt{fswebcam} capture, or an immediate image URL when an optional HTTP streamer is reachable.

\begin{figure}[H]
\centering
\begin{tikzpicture}[font=\scriptsize, box/.style={draw, rounded corners, align=center, minimum width=2.1cm, minimum height=0.65cm}, arr/.style={-{Stealth}, thick}]
\node[box, fill=leafgreen!14] (panel) {Pi panel\\Take Picture};
\node[box, right=0.9cm of panel, fill=blue!8] (api) {/api/devices};
\node[box, right=0.9cm of api, fill=teal!12, dashed] (hw) {Klipper\\/ agent};
\draw[arr] (panel) -- (api);
\draw[arr, dashed] (api) -- (hw);
\end{tikzpicture}
\caption{Tray Pi panel: Take Picture and Klipper URL management via devices API.}
\label{fig:pres-pi-panel}
\end{figure}

\subsubsection{Assumptions}
\begin{itemize}
  \item Offline empty states are shown when no device is linked or heartbeat is stale.
  \item Klipper/streamer URL updates validate http/https schemes before save.
\end{itemize}

\subsubsection{Responsibilities}
\begin{itemize}
  \item Trigger Take Picture and display returned capture or queue status.
  \item Persist Klipper URL edits (\texttt{updateKlipperUrl}; alias \texttt{updateMoonrakerUrl}).
  \item Support pose generation entry points for multi-angle capture.
\end{itemize}

\subsubsection{Subsystem interfaces}
\begin{table}[H]
\centering
\small
\begin{tabular}{|p{0.9cm}|p{3.4cm}|p{3.0cm}|p{3.5cm}|}
\hline
\textbf{ID} & \textbf{Description} & \textbf{Inputs} & \textbf{Outputs} \\
\hline
P-PI1 & Pi panel $\rightarrow$ Application & \texttt{capture} action & Image URL or queued command \\
\hline
P-PI2 & Pi panel $\rightarrow$ Application & Klipper URL & Updated \texttt{edge\_devices} row \\
\hline
P-PI3 & Pi panel $\rightarrow$ Application & Pose generate request & Pose sequence DTO \\
\hline
\end{tabular}
\caption{Tray Raspberry Pi panel and Take Picture subsystem interfaces}
\end{table}
